\subsection{The theory of Frédéric Riesz}

Let E be a separated locally convex space and let h be a compact endomorphism of E. Let \(\lambda \in K\) . The endomorphism \(\lambda h\) is compact, hence \(1_{E} - \lambda h\) is a Riesz endomorphism of E (cor. 2 of prop. 2 of III, p. 75); denote by \(N_{\lambda}\) and \(I_{\lambda}\) its nilspace and its conilspace. By prop. 6 of III, p. 47, we have the following properties:

\begin{enumerate}
\def\labelenumi{(\roman{enumi})}
\item
  The vector subspaces \(I_{\lambda}\) and \(N_{\lambda}\) of E are closed and stable under h;
\item
  The locally convex space E is the topological direct sum of \(I_{\lambda}\) and \(N_{\lambda}\) ;
\item
  The map \(1_{\mathrm{E}} - \lambda h\) induces by restriction an automorphism of \(\mathrm{I}_{\lambda}\) ;
\item
  The vector space \(N_{\lambda}\) is of finite dimension, and there exists an integer \(n_{\lambda} \geqslant 0\) such that \((1_{\mathrm{E}} - \lambda h)^{n_{\lambda}}(x) = 0\) for all \(x \in N_{\lambda}\) .
\end{enumerate}

These properties determine the spaces \(I_{\lambda}\) and \(N_{\lambda}\) in a unique manner (A, VIII, p. 26, remark 2).

Lemma 3. --- Let X be a topological space whose topology has a countable base. Every discrete subset of X is countable.

Let \(\mathcal{B}\) be a countable base of the topology of X and let D be a discrete subset of X. For each element \(x\) of D, there exists an element \(\mathrm{B}_x\) of \(\mathcal{B}\) such that \(\mathrm{B}_x \cap \mathrm{D} = \{x\}\) . The map from D to \(\mathcal{B}\) defined by \(x \mapsto \mathrm{B}_x\) is injective, whence the lemma.

THEOREM 4. --- Let E be a separated locally convex space and let h be a compact endomorphism of E. The set \(\Sigma\) of \(\lambda \in K\) such that \(1_{\mathrm{E}} - \lambda h\) is not an automorphism of E is a countable, closed, and discrete subset of the field K.

The set \(\Sigma\) is also the set of \(\lambda\in\mathrm{K}\) such that \(1_{\mathrm{E}}-\lambda h\) is not injective and the set of \(\lambda\in\mathrm{K}\) such that \(1_{\mathrm{E}}-\lambda h\) is not surjective.

There exists a Banach space G, a continuous linear map \(p: E \rightarrow G\) and a compact linear map \(q: G \rightarrow E\) such that \(h = q \circ p\) (III, p. 5, prop. 5). The endomorphism \(h' = p \circ q\) of G is compact, and \(1_{E} - \lambda h\) is an automorphism of E if and only if \(1_{G} - \lambda h'\) is an automorphism of G (prop. 10 of III, p. 49, e)). It therefore suffices for us to prove the theorem when E is a Banach space, which we henceforth assume.

Let \(\lambda \in K\) . Since the index of a Riesz endomorphism is zero, it is equivalent to say that the map \(1_{\mathrm{E}} - \lambda h\) is an automorphism of E, or that it is injective, or again that it is surjective.

The set \(\Sigma\) is closed in K since the set of automorphisms of E is open in \(\mathcal{L}(\mathrm{E})\) . Let \(\lambda\) be an element of \(\Sigma\) . We have \(\lambda \neq 0\) . Let N denote the nilspace of h and I its conilspace, and denote by \(h_{I}\) and \(h_{N}\) the endomorphisms of I and of N induced by h on these subspaces. Then \(1_{I} - \lambda h_{I}\) is an automorphism of I, and there exists a neighborhood U of \(\lambda\) in K such that \(1_{I} - \mu h_{I}\) is an automorphism of I for all \(\mu \in U\) . The endomorphism \(1_{N} - \lambda h_{N}\) of N is nilpotent; for all \(\mu \neq \lambda\) , we have

\[
1 _ {\mathrm{N}} - \mu h _ {\mathrm{N}} = \frac {\lambda - \mu}{\lambda} \Big (1 _ {\mathrm{N}} + \frac {\mu}{\lambda - \mu} (1 _ {\mathrm{N}} - \lambda h _ {\mathrm{N}}) \Big),
\]

hence \(1_{N} - \mu h_{N}\) is an automorphism of N. Consequently, the set \(U \cap \Sigma\) is reduced to the single element \(\lambda\) and the set \(\Sigma\) is discrete. It is countable by Lemma 3.

